\scanpage{61}
\mspage{54}

adjoint à gauche de

\[ u^*:D^-(I,A)\longrightarrow D^-(I'',A). \]

Mais je \deleted{\unsure{ne}} ne sais plus trop s'il faut s'attendre à pouvoir définir $u_!$ sur $D(I',A)$ \deleted{tout} entier, \deleted{\unsure{au}} \deleted{\unsure{en}} dehors des hypothèses genre dim. coh. finie. La triste vérité, c'est que j'ai $\pm$ oublié la théorie des foncteurs dérivés ! (Il faudrait que je demande \unsure{encore} à Jean Malgoire de m'envoyer des textes canoniques.)

\uline{Scholie} Si $(M,W)$ est quillénisable sur \deleted{\unsure{$\mathrm{Cat}$}} $\mathrm{Diag}\subset\mathrm{Cat}$, \inserted{et $=$ sans qu'il soit quillénisable,} il définit un prédérivateur \deleted{sur} $\mathbb{D}_{(M,W)}$ sur $\mathrm{Diag}$, \deleted{\unsure{qui est}} On \unsure{posera|écrira} parfois

\[ \mathrm{HO}_{M,W}(I)\ \text{ou}\ \mathrm{HO}_W(I)=\mathbb{D}_{M,W}(I). \]

L'hypothèse quillénisable implique, sur \inserted{\unsure{du moins pour|de même que}} ces ``catégories homotopiques, la structure \unsure{(relativement)} \unsure{``faible''} ``triangulée'', provenant des suites ``\uline{exactes}'' de suspension et de cosuspension\textsuperscript{*}. (La condition Quillénisable \unsure{assure} l'exactitude de ces suites \ldots)

% slanted

\margin{dans le cas \unsure{(particulier)|(cf. chapitre)} où $M$ est \unsure{pointée}, \unsure{sinon}, \unsure{par hypoth.|pour objet} \unsure{d'additivité} \illegible{1}, il y a \unsure{la structure|les structures} \unsure{associée|associées} aux \deleted{\unsure{foncteurs|carrés}} carrés cartésiens \unsure{resp.|et} cocartésiens. \unsure{Il suffit d'|Il s'agit d'}\illegible{1} \unsure{les carrés|les images} \unsure{et antéimages|et anti-images} \illegible{1} \unsure{caractérisent|caractérisés} \unsure{un autre type} \unsure{supplémentaire|propre} \unsure{de structure|à ces structures}.}

Les foncteurs $u^*:\mathrm{HO}_W(I)\unsure{-|\to}\mathrm{HO}_W(I')$
% archive sheet number: 60
